\documentclass[12pt]{article}
\usepackage[UTF8]{inputenc}
\usepackage{thaisupp}
\usepackage{enumitem}
\begin{document}
\title{แม่เหล็ก}
\maketitle
\section*{In-class Questions}
\begin{enumerate}[label=\arabic*.]
\item สนามแม่เหล็ก คือข้อใด \hfill \textbf{\small (Main)}
\begin{enumerate}[label=)]
\item ก) บริเวณรอบแม่เหล็กที่มีแรงแม่เหล็กทำให้วัสดุแม่เหล็กหรือสารแม่เหล็กเคลื่อนที่ได้
\item ข) บริเวณรอบแม่เหล็กที่มีแรงโน้มถ่วงทำให้วัสดุทุกชนิดเคลื่อนที่ได้
\item ค) บริเวณรอบแม่เหล็กที่มีสนามไฟฟ้าทำให้ประจุไฟฟ้าเคลื่อนที่ได้
\item ง) บริเวณรอบแม่เหล็กที่มีแรงเสียดทานทำให้วัสดุเคลื่อนที่ได้
\end{enumerate}
\item ขดลวดที่มีกระแสไฟฟ้าลัดผ่าน จะเกิดสนามแม่เหล็กรอบขดลวดนั้น ข้อใดถูกต้องเกี่ยวกับทิศทางของสนามแม่เหล็ก \hfill \textbf{\small (Main)}
\begin{enumerate}[label=)]
\item ก) ทิศทางเข้าขดลวดเป็นขั้วใต้ ทิศทางออกจากขดลวดเป็นขั้วเหนือ
\item ข) ทิศทางเข้าขดลวดเป็นขั้วเหนือ ทิศทางออกจากขดลวดเป็นขั้วใต้
\item ค) ทิศทางของสนามแม่เหล็กขึ้นกับทิศทางของกระแสไฟฟ้าเสมอ
\item ง) ทิศทางของสนามแม่เหล็กไม่ขึ้นกับทิศทางของกระแสไฟฟ้า
\end{enumerate}
\item เส้นแรงแม่เหล็ก มีคุณสมบัติอย่างไร \hfill \textbf{\small (Main)}
\begin{enumerate}[label=)]
\item ก) เส้นแรงแม่เหล็กเป็นเส้นตรงเสมอ
\item ข) เส้นแรงแม่เหล็กออกจากขั้วเหนือไปยังขั้วใต้ภายนอกแม่เหล็ก และจากขั้วใต้ไปยังขั้วเหนือภายในแม่เหล็ก
\item ค) เส้นแรงแม่เหล็กตัดกันเสมอ
\item ง) เส้นแรงแม่เหล็กไม่มีทิศทาง
\end{enumerate}
\item บริเวณใดมีสนามแม่เหล็กมากที่สุด \hfill \textbf{\small (Main)}
\begin{enumerate}[label=)]
\item ก) บริเวณกลางระหว่างขั้วทั้งสอง
\item ข) บริเวณใกล้ขั้วแม่เหล็กทั้งสอง
\item ค) บริเวณไกลจากขั้วแม่เหล็ก
\item ง) บริเวณรอบแม่เหล็กทุกจุดมีสนามแม่เหล็กเท่ากัน
\end{enumerate}
\item แม่เหล็กถาวรมีขั้วแม่เหล็กกี่ขั้ว และขั้วที่ชอบชี้ไปทางทิศเหนือทางภูมิศาสตร์เรียกว่าขั้วอะไร \hfill \textbf{\small (Main)}
\begin{enumerate}[label=)]
\item ก) 2 ขั้ว เรียกว่าขั้วใต้แม่เหล็ก
\item ข) 2 ขั้ว เรียกว่าขั้วเหนือแม่เหล็ก
\item ค) 4 ขั้ว เรียกว่าขั้วเหนือและขั้วใต้
\item ง) 2 ขั้ว เรียกว่าขั้วบวกและขั้วลบ
\end{enumerate}
\item เมื่อนำขั้วเหนือ (N) ของแม่เหล็ก两根ขั้วเข้าหากัน จะเกิดอะไรขึ้น \hfill \textbf{\small (Main)}
\begin{enumerate}[label=)]
\item ก) เกิดแรงดึงดูด
\item ข) เกิดแรงผลัก
\item ค) ไม่เกิดแรงใดๆ
\item ง) เกิดทั้งแรงดึงดูดและแรงผลัก
\end{enumerate}
\item เมื่อนำขั้วเหนือ (N) และขั้วใต้ (S) ของแม่เหล็ก两根ขั้วเข้าหากัน จะเกิดอะไรขึ้น \hfill \textbf{\small (Main)}
\begin{enumerate}[label=)]
\item ก) เกิดแรงผลัก
\item ข) เกิดแรงดึงดูด
\item ค) ไม่เกิดแรงใดๆ
\item ง) แม่เหล็กจะหมุนตามแนวตั้ง
\end{enumerate}
\item สนามแม่เหล็กโลกเกิดจากอะไร \hfill \textbf{\small (Main)}
\begin{enumerate}[label=)]
\item ก) แกนเหล็กในโลก
\item ข) กระแสไฟฟ้าที่ไหลวนในแก่นโลกซึ่งเป็นโลหะหลอมละลาย
\item ค) แม่เหล็กถาวรขนาดใหญ่ใต้ผิวโลก
\item ง) การเคลื่อนที่ของอิเล็กตรอนรอบนิวเคลียส
\end{enumerate}
\item ข้อใดเกี่ยวกับสนามแม่เหล็กโลกไม่ถูกต้อง \hfill \textbf{\small (Main)}
\begin{enumerate}[label=)]
\item ก) ขั้วเหนือแม่เหล็กโลกอยู่ใกล้ขั้วใต้ทางภูมิศาสตร์
\item ข) ขั้วใต้แม่เหล็กโลกอยู่ใกล้ขั้วเหนือทางภูมิศาสตร์
\item ค) ทิศทางของสนามแม่เหล็กโลกวัดได้โดยเข็มทิศ
\item ง) ขั้วเหนือแม่เหล็กโลกอยู่ที่ขั้วเหนือทางภูมิศาสตร์พอดี
\end{enumerate}
\item เข็มทิศชี้ไปทางทิศเหนือเสมอ เพราะเหตุใด \hfill \textbf{\small (Main)}
\begin{enumerate}[label=)]
\item ก) โลกหมุนรอบตัวเอง
\item ข) ขั้วเหนือของแม่เหล็กในเข็มทิศถูกดึงดูดไปยังขั้วใต้แม่เหล็กโลก
\item ค) แรงโน้มถ่วงของโลกดึงเข็มทิศ
\item ง) สนามไฟฟ้าของโลกผลักเข็มทิศ
\end{enumerate}
\item เส้นแรงแม่เหล็กจากขั้วหนึ่งไปยังอีกขั้วหนึ่ง ข้อใดไม่ใช่คุณสมบัติของเส้นแรงแม่เหล็ก \hfill \textbf{\small (Main)}
\begin{enumerate}[label=)]
\item ก) เส้นแรงแม่เหล็กจะไม่ตัดกัน
\item ข) เส้นแรงแม่เหล็กจะพยายามหดสั้นลงเสมอ
\item ค) เส้นแรงแม่เหล็กจะเข้าหากันเสมอ
\item ง) ความหนาแน่นของเส้นแรงบ่งบอกความเข้มของสนามแม่เหล็ก
\end{enumerate}
\item วัสดุในข้อใดต่อไปนี้ที่สามารถทำให้เกิดสนามแม่เหล็กได้ \hfill \textbf{\small (Main)}
\begin{enumerate}[label=)]
\item ก) เหล็ก อะลูมิเนียม และทองแดง
\item ข) เหล็ก นิกเกิล และโคบอลต์
\item ค) อะลูมิเนียม ตะกั่ว และสังกะสี
\item ง) ทองแดง สนิมเหล็ก และเงิน
\end{enumerate}
\item สนามแม่เหล็ก B ที่ระยะ r จากขดลวดยาว L ที่มีกระแส I ไหลผ่าน ขึ้นอยู่กับปริมาณใดเป็นหลัก \hfill \textbf{\small (Main)}
\begin{enumerate}[label=)]
\item ก) ความต้านทานของขดลวด
\item ข) จำนวนรอบของขดลวดและกระแสไฟฟ้า
\item ค) แรงดันไฟฟ้าของแหล่งจ่าย
\item ง) ความยาวของขดลวดเท่านั้น
\end{enumerate}
\item เมื่อนำแผ่นกระดาษมาวางทับแท่งแม่เหล็กแล้วโรยผงเหล็กบนแผ่นกระดาษ ผงเหล็กจะจัดตัวเป็นรูปแบบใด \hfill \textbf{\small (Main)}
\begin{enumerate}[label=)]
\item ก) เป็นเส้นตรงขนานกันไปทุกทิศทาง
\item ข) เป็นเส้นโค้งเชื่อมจากขั้วเหนือไปขั้วใต้
\item ค) เป็นวงกลมรอบแม่เหล็ก
\item ง) กระจายตัวสม่ำเสมอทั่วแผ่นกระดาษ
\end{enumerate}
\item มุมเบี้ยว (declination) ของเข็มทิศคืออะไร \hfill \textbf{\small (Main)}
\begin{enumerate}[label=)]
\item ก) มุมที่เข็มทิศทำกับแนวระดับ
\item ข) มุมที่ขั้วเหนือแม่เหล็กทำกับขั้วเหนือทางภูมิศาสตร์
\item ค) มุมที่เข็มทิศทำกับแนวดิ่ง
\item ง) มุมระหว่างสนามแม่เหล็กกับสนามไฟฟ้า
\end{enumerate}
\item หน่วยของสนามแม่เหล็ก B ในระบบ SI คืออะไร \hfill \textbf{\small (Main)}
\begin{enumerate}[label=)]
\item ก) เทสลา (Tesla, T)
\item ข) เวเบอร์ (Weber, Wb)
\item ค) แอมแปร์ต่อเมตร (A/m)
\item ง) โอห์ม (Ω)
\end{enumerate}
\item ถ้าตัดแม่เหล็กแท่งออกเป็นสองท่อน แต่ละท่อนจะมีขั้วแม่เหล็กอย่างไร \hfill \textbf{\small (Main)}
\begin{enumerate}[label=)]
\item ก) ท่อนหนึ่งมีขั้วเหนือ อีกท่อนมีขั้วใต้
\item ข) ทั้งสองท่อนมีขั้วเหนือและขั้วใต้ครบ
\item ค) ท่อนหนึ่งมีขั้วเหนือ อีกท่อนมีขั้วใต้ และไม่มีขั้วอีกขั้ว
\item ง) ทั้งสองท่อนไม่มีขั้วแม่เหล็ก
\end{enumerate}
\item ขดลวดรูปวงแหวนที่มีรัศมี r มีกระแส I ไหลผ่าน สนามแม่เหล็กที่จุดศูนย์กลางของขดลวดวงแหวนมีค่าเท่าใด \hfill \textbf{\small (Main)}
\begin{enumerate}[label=)]
\item ก) B = (μ₀I)/(2r)
\item ข) B = (μ₀I)/(2πr)
\item ค) B = (μ₀NI)/(2r)
\item ง) B = (μ₀I)/(r²)
\end{enumerate}
\item มุมเอียง (inclination) ของสนามแม่เหล็กโลกคืออะไร \hfill \textbf{\small (Main)}
\begin{enumerate}[label=)]
\item ก) มุมระหว่างสนามแม่เหล็กโลกกับแนวระดับ
\item ข) มุมระหว่างสนามแม่เหล็กโลกกับแนวดิ่ง
\item ค) มุมระหว่างแกนโลกกับแกนแม่เหล็กโลก
\item ง) มุมระหว่างเข็มทิศกับแนวเหนือ-ใต้
\end{enumerate}
\item ฟลักซ์แม่เหล็ก (magnetic flux) คืออะไร \hfill \textbf{\small (Main)}
\begin{enumerate}[label=)]
\item ก) ผลคูณระหว่างสนามแม่เหล็กกับพื้นที่
\item ข) ผลคูณระหว่างสนามแม่เหล็กกับมุม
\item ค) ผลคูณระหว่างกระแสไฟฟ้ากับเวลา
\item ง) ผลคูณระหว่างแรงกับระยะทาง
\end{enumerate}
\item สนามแม่เหล็กที่จุดศูนย์กลางของขดลวดโซลินอยด์ที่ยาวมาก มีค่าเท่าใด \hfill \textbf{\small (Main)}
\begin{enumerate}[label=)]
\item ก) B = μ₀NI/L
\item ข) B = μ₀N²IL
\item ค) B = μ₀I/(2πr)
\item ง) B = μ₀IL/(2r)
\end{enumerate}
\item ถ้าเพิ่มระยะห่างระหว่างขั้วแม่เหล็กต่างชนิดสองขั้วเป็น 2 เท่า แรงดึงดูดจะเปลี่ยนแปลงอย่างไร \hfill \textbf{\small (Main)}
\begin{enumerate}[label=)]
\item ก) เพิ่มเป็น 2 เท่า
\item ข) ลดลงเป็น 1/2 เท่า
\item ค) ลดลงเป็น 1/4 เท่า
\item ง) ไม่เปลี่ยนแปลง
\end{enumerate}
\item ข้อใดไม่ใช่วิธีการทำให้วัสดุเฟอร์โรแมกเนติกกลายเป็นแม่เหล็กถาวร \hfill \textbf{\small (Main)}
\begin{enumerate}[label=)]
\item ก) นำไปวางในสนามแม่เหล็กที่แรงพอสมควร
\item ข) นำไปเสียบเข้ากับวงจรไฟฟ้า
\item ง) นำไปทำให้ร้อนแล้วทิ้งให้เย็นลงในสนามแม่เหล็ก
\end{enumerate}
\item ถ้าวางเข็มทิศไว้ใกล้สายไฟที่มีกระแสไฟฟ้าคงที่ เข็มทิศจะเบนออกจากแนวเหนือ-ใต้ แสดงว่า \hfill \textbf{\small (Main)}
\begin{enumerate}[label=)]
\item ก) สนามไฟฟ้าทำให้เข็มเบน
\item ข) สนามแม่เหล็กที่เกิดจากกระแสไฟฟ้าทำให้เข็มเบน
\item ค) ความร้อนจากสายไฟทำให้เข็มเบน
\item ง) แรงดึงดูดระหว่างกระแสกับเข็มทิศทำให้เบน
\end{enumerate}
\item ขดลวดตรงยาวมากมีจำนวนรอบต่อหนึ่งหน่วยความยาว n = 1000 รอบ/เมตร มีกระแส I = 2 A ไหลผ่าน สนามแม่เหล็กภายในขดลวดมีค่าเท่าใด (μ₀ = 4π × 10⁻⁷ T·m/A) \hfill \textbf{\small (Main)}
\begin{enumerate}[label=)]
\item ก) 2.5 × 10⁻³ T
\item ข) 2.5 × 10⁻⁴ T
\item ค) 2.5 × 10⁻² T
\item ง) 25 × 10⁻⁴ T
\end{enumerate}
\item สนามแม่เหล็กเป็นปริมาณเวกเตอร์หรือปริมาณสเกลาร์ \hfill \textbf{\small (Extra)}
\begin{enumerate}[label=)]
\item ก) ปริมาณสเกลาร์
\item ข) ปริมาณเวกเตอร์
\item ค) ขึ้นกับกรณี
\item ง) เป็นได้ทั้งสองอย่าง
\end{enumerate}
\item ข้อใดไม่ใช่คุณสมบัติของแม่เหล็ก \hfill \textbf{\small (Extra)}
\begin{enumerate}[label=)]
\item ก) แม่เหล็กมีสองขั้ว คือขั้วเหนือและขั้วใต้
\item ข) ขั้วเดียวกันจะผลักกัน
\item ค) ขั้วต่างกันจะดึงดูดกัน
\item ง) สามารถแยกขั้วแม่เหล็กออกจากกันได้โดยการตัด
\end{enumerate}
\item เส้นแรงแม่เหล็กจะหนาแน่นมากที่สุดบริเวณใด \hfill \textbf{\small (Extra)}
\begin{enumerate}[label=)]
\item ก) บริเวณกลางระหว่างขั้วทั้งสอง
\item ข) บริเวณขั้วแม่เหล็ก
\item ค) บริเวณไกลจากขั้วแม่เหล็ก
\item ง) บริเวณรอบกลางแม่เหล็ก
\end{enumerate}
\item ข้อใดไม่ถูกต้องเกี่ยวกับสนามแม่เหล็กโลก \hfill \textbf{\small (Extra)}
\begin{enumerate}[label=)]
\item ก) มีลักษณะเหมือนมีแท่งแม่เหล็กขนาดใหญ่วางอยู่ในแก่นโลก
\item ข) ขั้วเหนือแม่เหล็กโลกอยู่ที่ขั้วเหนือทางภูมิศาสตร์พอดี
\item ค) ทำให้เข็มทิศชี้ไปทางทิศเหนือได้
\item ง) เกิดจากกระแสไฟฟ้าที่ไหลวนในแก่นโลก
\end{enumerate}
\item สารในข้อใดทั้งหมดเป็นสารแม่เหล็ก \hfill \textbf{\small (Extra)}
\begin{enumerate}[label=)]
\item ก) เหล็ก อะลูมิเนียม ทองแดง
\item ข) เหล็ก นิกเกิล โคบอลต์
\item ค) นิกเกิล ตะกั่ว สังกะสี
\item ง) โคบอลต์ ดีบุก แมกนีเซียม
\end{enumerate}
\item แรงระหว่างขั้วแม่เหล็กสองขั้วที่ห่างกันระยะ r แปรผกผันกับข้อใด \hfill \textbf{\small (Extra)}
\begin{enumerate}[label=)]
\item ก) r
\item ข) r²
\item ค) 1/r
\item ง) 1/r²
\end{enumerate}
\item เมื่อเพิ่มกระแสไฟฟ้าในขดลวดเป็น 2 เท่า สนามแม่เหล็กที่เกิดจะเปลี่ยนแปลงอย่างไร \hfill \textbf{\small (Extra)}
\begin{enumerate}[label=)]
\item ก) เพิ่มเป็น 2 เท่า
\item ข) เพิ่มเป็น 4 เท่า
\item ค) ลดลงเป็น 1/2 เท่า
\item ง) ไม่เปลี่ยนแปลง
\end{enumerate}
\item เส้นแรงแม่เหล็กมีคุณสมบัติข้อใดที่ไม่ถูกต้อง \hfill \textbf{\small (Extra)}
\begin{enumerate}[label=)]
\item ก) เส้นแรงแม่เหล็กไม่ตัดกัน
\item ข) เส้นแรงแม่เหล็กเป็นเส้นปิด
\item ค) เส้นแรงแม่เหล็กจะพยายามหดสั้นลงเสมอ
\item ง) เส้นแรงแม่เหล็กไปจากขั้วใต้ไปขั้วเหนือเสมอ
\end{enumerate}
\item มุมเบี้ยว (declination) มีค่าเป็นบวกเมื่อใด \hfill \textbf{\small (Extra)}
\begin{enumerate}[label=)]
\item ก) เมื่อขั้วเหนือแม่เหล็กชี้ไปทางตะวันตกของทิศเหนือทางภูมิศาสตร์
\item ข) เมื่อขั้วเหนือแม่เหล็กชี้ไปทางตะวันออกของทิศเหนือทางภูมิศาสตร์
\item ค) เมื่อขั้วเหนือแม่เหล็กชี้ตรงทิศเหนือทางภูมิศาสตร์
\item ง) เมื่อเข็มทิศอยู่ในแนวระดับ
\end{enumerate}
\item ขดลวดทรงกลม (toroid) มีสนามแม่เหล็กอยู่ที่ใด \hfill \textbf{\small (Extra)}
\begin{enumerate}[label=)]
\item ก) ภายในขดลวดเท่านั้น
\item ข) ภายนอกขดลวดเท่านั้น
\item ค) ทั้งภายในและภายนอกขดลวด
\item ง) ไม่มีสนามแม่เหล็กทั้งภายในและภายนอก
\end{enumerate}
\item สนามแม่เหล็กรอบสายไฟตรงที่มีกระแสไฟฟ้าไหล มีลักษณะอย่างไร \hfill \textbf{\small (Extra)}
\begin{enumerate}[label=)]
\item ก) เป็นเส้นตรงขนานกับสายไฟ
\item ข) เป็นวงกลมรอบสายไฟ
\item ค) เป็นเส้นตรงออกจากสายไฟ
\item ง) เป็นเส้นโค้งเข้าหาสายไฟ
\end{enumerate}
\item สนามแม่เหล็กภายในขดลวดโซลินอยด์ที่ยาวมากจะมีลักษณะอย่างไร \hfill \textbf{\small (Extra)}
\begin{enumerate}[label=)]
\item ก) มีลักษณะเป็นเส้นโค้งวงปิด
\item ข) มีลักษณะเป็นเส้นขนานตรง หนาแน่นสม่ำเสมอ
\item ค) มีลักษณะกระจายออกจากกัน
\item ง) มีลักษณะเป็นวงกลมรอบขดลวด
\end{enumerate}
\item เมื่อนำแม่เหล็กให้สัมผัสกับเหล็กไร้ขั้ว แม่เหล็กจะทำให้เหล็กไร้ขั้วเกิดอะไรขึ้น \hfill \textbf{\small (Extra)}
\begin{enumerate}[label=)]
\item ก) เหล็กไร้ขั้วจะกลายเป็นแม่เหล็กถาวร
\item ข) เหล็กไร้ขั้วจะถูกเหนี่ยวนำให้มีขั้วตรงข้ามกับแม่เหล็กที่สัมผัส
\item ค) เหล็กไร้ขั้วจะถูกเหนี่ยวนำให้มีขั้วเดียวกับแม่เหล็กที่สัมผัส
\item ง) เหล็กไร้ขั้วจะไม่เปลี่ยนแปลง
\end{enumerate}
\item มุมเอียง (inclination) ที่ขั้วแม่เหล็กโลกมีค่าเท่าใด \hfill \textbf{\small (Extra)}
\begin{enumerate}[label=)]
\item ก) 0 องศา
\item ข) 45 องศา
\item ค) 90 องศา
\item ง) 180 องศา
\end{enumerate}
\item ถ้าใช้เข็มทิศวัดทิศทางสนามแม่เหล็ก ข้อใดถูกต้องเกี่ยวกับทิศทางของสนามแม่เหล็ก \hfill \textbf{\small (Extra)}
\begin{enumerate}[label=)]
\item ก) ทิศทางที่ขั้วเหนือของเข็มทิศชี้ คือทิศทางของสนามแม่เหล็ก
\item ข) ทิศทางที่ขั้วใต้ของเข็มทิศชี้ คือทิศทางของสนามแม่เหล็ก
\item ค) ทิศทางของสนามแม่เหล็กตั้งฉากกับทิศทางที่เข็มทิศชี้
\item ง) ทิศทางของสนามแม่เหล็กไม่เกี่ยวข้องกับทิศทางเข็มทิศ
\end{enumerate}
\item ขดลวดวงแหวนรัศมี 0.1 เมตร มีกระแส 5 A ไหลผ่าน สนามแม่เหล็กที่จุดศูนย์กลางของขดลวดมีค่าเท่าใด (μ₀ = 4π × 10⁻⁷ T·m/A) \hfill \textbf{\small (Extra)}
\begin{enumerate}[label=)]
\item ก) 3.14 × 10⁻⁵ T
\item ข) 3.14 × 10⁻⁴ T
\item ค) 3.14 × 10⁻³ T
\item ง) 3.14 × 10⁻⁶ T
\end{enumerate}
\item ขดลวดโซลินอยด์ยาว 0.5 เมตร มีจำนวนรอบ 1000 รอบ มีกระแส 4 A ไหลผ่าน สนามแม่เหล็กภายในโซลินอยด์มีค่าเท่าใด \hfill \textbf{\small (Extra)}
\begin{enumerate}[label=)]
\item ก) 1.0 × 10⁻² T
\item ข) 1.0 × 10⁻³ T
\item ค) 1.0 × 10⁻⁴ T
\item ง) 1.0 × 10⁻¹ T
\end{enumerate}
\item ถ้าแรงระหว่างขั้วแม่เหล็กสองขั้วที่ห่างกัน 0.1 เมตร มีค่า 4 N เมื่อเพิ่มระยะห่างเป็น 0.2 เมตร แรงจะมีค่าเท่าใด \hfill \textbf{\small (Extra)}
\begin{enumerate}[label=)]
\item ก) 2 N
\item ข) 1 N
\item ค) 0.5 N
\item ง) 8 N
\end{enumerate}
\item สายไฟตรงยาวมากมีกระแส 10 A ไหลผ่าน สนามแม่เหล็กที่ระยะ 0.05 เมตรจากสายไฟมีค่าเท่าใด \hfill \textbf{\small (Extra)}
\begin{enumerate}[label=)]
\item ก) 4 × 10⁻⁵ T
\item ข) 4 × 10⁻⁴ T
\item ค) 2 × 10⁻⁴ T
\item ง) 2 × 10⁻⁵ T
\end{enumerate}
\item ฟลักซ์แม่เหล็กที่ผ่านพื้นที่ 0.02 ตารางเมตร ที่ทำมุม 60° กับสนามแม่เหล็ก 5 × 10⁻³ T มีค่าเท่าใด \hfill \textbf{\small (Extra)}
\begin{enumerate}[label=)]
\item ก) 5 × 10⁻⁵ Wb
\item ข) 5 × 10⁻⁴ Wb
\item ค) 5 × 10⁻³ Wb
\item ง) 5 × 10⁻⁶ Wb
\end{enumerate}
\item ที่เส้นศูนย์สูตร มุมเอียงของสนามแม่เหล็กโลกมีค่าเท่าใด \hfill \textbf{\small (Extra)}
\begin{enumerate}[label=)]
\item ก) 0°
\item ข) 45°
\item ค) 90°
\item ง) 180°
\end{enumerate}
\item แม่เหล็กไฟฟ้าทำงานต่างจากแม่เหล็กถาวรอย่างไร \hfill \textbf{\small (Extra)}
\begin{enumerate}[label=)]
\item ก) แม่เหล็กไฟฟ้ามีขั้วแม่เหล็ก 2 ขั้ว ส่วนแม่เหล็กถาวรมี 4 ขั้ว
\item ข) แม่เหล็กไฟฟ้าสามารถควบคุมความเข้มสนามได้โดยการเปลี่ยนกระแส
\item ค) แม่เหล็กถาวรมีสนามแม่เหล็กคงที่เสมอ
\item ง) แม่เหล็กไฟฟ้าทำจากเหล็กกล้า ส่วนแม่เหล็กถาวรทำจากเหล็กอ่อน
\end{enumerate}
\item ขดลวดตรงยาวมากมีจำนวนรอบต่อหนึ่งหน่วยความยาว 2000 รอบ/เมตร มีกระแส 3 A สนามแม่เหล็กภายในขดลวดมีค่าเท่าใด \hfill \textbf{\small (Extra)}
\begin{enumerate}[label=)]
\item ก) 7.54 × 10⁻³ T
\item ข) 7.54 × 10⁻⁴ T
\item ค) 7.54 × 10⁻² T
\item ง) 7.54 × 10⁻⁵ T
\end{enumerate}
\item สนามแม่เหล็ก B = 0.1 T ทำมุม 30° กับพื้นที่ A = 0.05 ตารางเมตร ฟลักซ์แม่เหล็กที่ผ่านพื้นที่นี้มีค่าเท่าใด \hfill \textbf{\small (Extra)}
\begin{enumerate}[label=)]
\item ก) 2.5 × 10⁻³ Wb
\item ข) 5.0 × 10⁻³ Wb
\item ค) 2.5 × 10⁻⁴ Wb
\item ง) 5.0 × 10⁻⁴ Wb
\end{enumerate}
\item ขั้วแม่เหล็กสองขั้วต่างชนิดมีความเข้มขั้ว m₁ และ m₂ ห่างกัน r แรงดึงดูดระหว่างขั้วทั้งสองมีค่าเท่าใด \hfill \textbf{\small (Extra)}
\begin{enumerate}[label=)]
\item ก) F = k·m₁m₂/r
\item ข) F = k·m₁m₂/r²
\item ค) F = k·m₁/r²
\item ง) F = k·m₁m₂·r²
\end{enumerate}
\item ขดลวดวงแหวน N รอบ รัศมี r มีกระแส I ไหลผ่าน ถ้าเพิ่มจำนวนรอบเป็น 2 เท่า และรัศมีเป็น 1/2 เท่า สนามแม่เหล็กที่จุดศูนย์กลางจะเปลี่ยนแปลงอย่างไร \hfill \textbf{\small (Extra)}
\begin{enumerate}[label=)]
\item ก) เพิ่มเป็น 2 เท่า
\item ข) เพิ่มเป็น 4 เท่า
\item ค) เพิ่มเป็น 8 เท่า
\item ง) เพิ่มเป็น 16 เท่า
\end{enumerate}
\item เหตุใดแม่เหล็กที่ทำจากเหล็กอ่อนจึงใช้ทำแกนของแม่เหล็กไฟฟ้าได้ดี \hfill \textbf{\small (Extra)}
\begin{enumerate}[label=)]
\item ก) เหล็กอ่อนมีความเข้มขั้วแม่เหล็กสูง
\item ข) เหล็กอ่อนสามารถเป็นแม่เหล็กได้ง่ายแต่คงสภาพแม่เหล็กได้ไม่นาน
\item ค) เหล็กอ่อนมีความต้านทานไฟฟ้าสูง
\item ง) เหล็กอ่อนนำความร้อนได้ดี
\end{enumerate}
\item ขดลวดโซลินอยด์ยาว 0.4 เมตร มีเส้นผ่านศูนย์กลาง 0.1 เมตร มีจำนวนรอบ 800 รอบ มีกระแส 2 A ไหลผ่าน สนามแม่เหล็กภายในโซลินอยด์มีค่าเท่าใด \hfill \textbf{\small (Extra)}
\begin{enumerate}[label=)]
\item ก) 5 × 10⁻⁴ T
\item ข) 5 × 10⁻³ T
\item ค) 1.6 × 10⁻³ T
\item ง) 1.6 × 10⁻⁴ T
\end{enumerate}
\item ถ้าสนามแม่เหล็กโลกที่ผิวโลกมีค่าประมาณ 5 × 10⁻⁵ T ที่ระยะ 1 เมตรจากขั้วแม่เหล็กของแท่งแม่เหล็กขนาดเล็กมีค่า 2 × 10⁻³ T ข้อใดถูกต้อง \hfill \textbf{\small (Extra)}
\begin{enumerate}[label=)]
\item ก) สนามแม่เหล็กของแท่งแม่เหล็กมากกว่าสนามแม่เหล็กโลก 40 เท่า
\item ข) สนามแม่เหล็กของแท่งแม่เหล็กน้อยกว่าสนามแม่เหล็กโลก 40 เท่า
\item ค) สนามแม่เหล็กทั้งสองมีค่าเท่ากัน
\item ง) สนามแม่เหล็กของแท่งแม่เหล็กมากกว่าสนามแม่เหล็กโลก 250 เท่า
\end{enumerate}
\item ขดลวดวงแหวน 50 รอบ มีรัศมี 0.2 เมตร มีกระแส 4 A ถ้านำขดลวดวงแหวนนี้ไปวางในสนามแม่เหล็กภายนอก 0.1 T ที่ทำมุม 30° กับแกนของขดลวด ฟลักซ์แม่เหล็กรวมที่ผ่านขดลวดมีค่าเท่าใด \hfill \textbf{\small (Extra)}
\begin{enumerate}[label=)]
\item ก) 3.46 × 10⁻³ Wb
\item ข) 1.73 × 10⁻³ Wb
\item ค) 6.93 × 10⁻³ Wb
\item ง) 8.66 × 10⁻³ Wb
\end{enumerate}
\end{enumerate}
\end{document}
